% ============================================================
% 嵌入式对接教程（一）· Part II：工程实现详解
% 由 emb1_tutorial.tex \input 引入（共用同一份导言区与主题色）
% ============================================================

\part{工程实现详解}

\begin{tipbox}{这一部分写什么}
把集成与仿真设施逐文件讲透：官方链路的接入细节、两个行为替身
（divider\_gate\_sim 52 行 / ph1p\_logic\_eram\_stub 160 行）、
system\_top 顶层、虚拟 RAW10 传感器与整链 TB，
以及 8 轮迭代里属于本册的 5 轮复盘。
\end{tipbox}

% ============================================================
\section{工程总览}

\subsection{本册涉及的文件}

\begin{center}
\begin{tabular}{llr l}
\toprule
\rowcolor{rowtint}\textbf{文件} & \textbf{类型} & \textbf{行数} & \textbf{职责} \\
\midrule
\code{rtl/system\_top.v} & RTL & 154 & 集成顶层（仿真/TD 双用） \\
\code{rtl/divider\_gate\_sim.v} & 替身 & 52 & 门级除法器的行为级替身 \\
\code{rtl/ph1p\_logic\_eram\_stub.v} & 替身 & 160 & TD ERAM 原语行为替身 \\
\code{tb/tb\_system\_integration.v} & 测试 & 163 & 整链对接四判据 \\
\code{sim/isp\_top.v} 等 & 官方副本 & --- & 官方明码 RTL（仿真用） \\
\code{sim/raw10\_beats.hex} & 向量 & --- & 4 帧虚拟传感器（230400 拍/帧） \\
\midrule
\code{tools/gen\_raw10\_sensor.py} & 脚本 & --- & 虚拟传感器生成器 \\
\bottomrule
\end{tabular}
\end{center}

\subsection{四步交付与分层证明}

嵌入式交付的四步（环境/驱动/配置/联调）在报告里各占一格；
本册覆盖第（1）步的仿真环境与第（4）步的集成联调，
第二册覆盖第（2）（3）步（驱动与摄像头配置）。分层证明的对应：

\begin{center}
\begin{tabular}{lll}
\toprule
\rowcolor{rowtint}\textbf{层} & \textbf{主张} & \textbf{证据} \\
\midrule
契约层 & 视觉兼容 96bit 契约 & tb\_vision\_top 14 项回归 \\
挂载层 & 官方 ISP + 视觉可共存于一个顶层 & elaboration 零错误 \\
行为层 & 整链数据流贯通（受 demosaic 限制） & tb\_system\_integration 判据 A/C \\
外设层 & UART 遥测逐字节正确 & tb\_uart\_telemetry（第二册详述） \\
\bottomrule
\end{tabular}
\end{center}

% ============================================================
\section{官方链路的接入细节}

\subsection{isp\_top 实例化的两个坑}

\begin{Verilog}
isp_top u_isp (
    ...
    .I_tvalid  (I_raw10_valid & isp_tready),  // 与上游就绪相与
    .O_tready  (isp_tready),                  // 它是输出: 对上游指示
    .I_tready  (isp_up_tready)                // 也是输出(命名反向)
);
\end{Verilog}

坑一：\code{I\_tready} 实际方向是输出，接常量直接编译报错；
坑二：\code{I\_tvalid} 要与 \code{isp\_tready} 相与——真实流控语义，
上游数据只有在 ISP 就绪时才有效。\textbf{读官方代码先读端口注释，
不要相信前缀命名}。

\subsection{RAW10 进、128bit 出}

数据传输的实际形状：

\begin{Verilog}
40bit/拍 (4 像素 x 10bit) x 320 拍/行 x 720 行 = 230400 拍/帧
   -> isp_top 内部: 解包/黑电平/去马赛克/白平衡 (整拍流水)
   -> 128bit/拍 (8 像素 x 24bit RGB) 输出
   -> data128_96 纯组合重排 -> 96bit (4 像素)
\end{Verilog}

% ============================================================
\section{仿真替身详解}

\subsection{divider\_gate\_sim：一个"假装很慢"的除法器}

真件是 7355 行 AL\_* 门级网表；替身用 FPGA 也不擅长、但仿真器
支持的行为级除法：

\begin{Verilog}
if (!busy && start && !done) begin
    num_q <= numerator; den_q <= denominator; busy <= 1; wait_cnt <= 0;
end else if (busy) begin
    if (wait_cnt == 6'd15) begin            // 16 拍"伪装"门级延迟
        if (den_q != 0) begin
            quotient  <= num_q / den_q;      // 行为级除法(仿真器算)
            remainder <= num_q % den_q;
        end else begin
            quotient  <= {48{1'b1}};         // 除零: 全 1
            remainder <= num_q;
        end
        done <= 1'b1; busy <= 1'b0;
    end else wait_cnt <= wait_cnt + 1;
end
\end{Verilog}

两个设计点：\textbf{16 拍延迟}是为了给下游状态机"有等待可等"
（AWB 的控制环本来就按若干拍设计）；\textbf{除零给全 1}是约定的
哨兵值。再次强调：TD 综合时这个文件\textbf{不进工程}，
真网表在位。

\subsection{ph1p\_logic\_eram\_stub：双口 RAM 行为模型}

官方 blk\_mem 生成的 RTL 例化 TD 原语，参数/端口极多。
替身的编写流程是"从实例化现场抄"：

\begin{center}
\begin{tabular}{ll}
\toprule
\rowcolor{rowtint}\textbf{项目} & \textbf{处理} \\
\midrule
参数（DATA\_WIDTH\_A 等） & 全部声明，多数不影响行为 \\
端口 clka/clkb/cea/ceb & 时钟与使能，行为使用 \\
端口 dia/dib/doa/dob & 数据 \\
端口 bea/beb & 字节使能（全字写时忽略） \\
\textbf{ECC 端口} & sbiterr/dbiterr 输出；*inj 输入 \\
\bottomrule
\end{tabular}
\end{center}

\begin{warnbox}{ECC 端口方向抄反过一次}
第一版替身按"错误信号应该是输入吧"想当然，把 \code{sbiterr/dbiterr}
写成输入——ModelSim 编译通过（方向错误在 Verilog 端口连接里
不总是报错），但上板/综合语义全乱。修复：回官方实例化清单
逐个核对方向。\textbf{替身的价值 = 接口逐位忠实，不是"能编过"}。
\end{warnbox}

% ============================================================
\section{system\_top 逐段}

\subsection{信号流与"三个域"}

\begin{Verilog}
I_clk (axi4s_video_aclk, 74.25M 级):  ISP / 视觉 / 遥测 / 声光
I_pwm_clk (50MHz):                    gimbal_ctrl / pwm_servo
两者在 vision_top 内部经 pulse_sync/cdc_4phase 跨域(视觉已备)
\end{Verilog}

顶层本身\textbf{不做任何跨域}——跨域全部封装在 vision\_top 内部，
顶层只负责摆位置、接按键、引输出。这是"分层"的又一处体现。

\subsection{按键 glue 与遥测的层次化引用}

\begin{Verilog}
if (key0_trig) preset_idx <= preset_idx + 3'd1;   // KEY0: 阈值组
if (key1_trig) track_en   <= ~track_en;           // KEY1: 跟踪使能

uart_telemetry u_uart (
    .I_frame_tick (u_vision.t_tick),   // 层次化引用: 顶层特权
    .I_target_valid(u_vision.t_valid),
    .I_cx (u_vision.t_cx), .I_cy (u_vision.t_cy), ...);
\end{Verilog}

遥测拿的是视觉内部信号（\code{t\_tick/t\_valid/t\_cx/t\_cy}）。
选层次化引用而不是加端口：\textbf{遥测是"观测者"，不该改变被观测
对象的接口}——视觉模块作为交付给上板的产物保持干净；
观测通道只在集成顶层存在。

\subsection{仿真与上板的双用语法}

顶层文件同时服务两个场景：ModelSim 仿真与 TD 对接。
它\textbf{不含}三种黑盒（DDR2 / MIPI-PHY / HDMI）——仿真里
以帧转储代替 HDMI，DDR 写流空置；上板时这些角色由官方
wrapper 承担，本文件的视觉挂载段整体搬过去。

% ============================================================
\section{测试平台设计}

\subsection{虚拟传感器生成器}

\begin{Verilog}
raw = 512 灰底 (10bit)
红块 80x80 @ (500,300): R 位点=900, G/B 位点=64
噪声: 0.2% 像素替换为随机值 (default_rng(500) 固定种子)
打包: 40bit = {P0[39:30], P1[29:20], P2[19:10], P3[9:0]}
标志: 每行第 0 拍 sof、第 319 拍 eol, 单独文件一行一值
\end{Verilog}

红块中心 (540, 340) 与视觉 TPG 测试的红块位置接近——
整链仿真复用视觉侧已有的期望坐标口径（±40px）。

\subsection{TB 的四个检查器}

\begin{itemize}
  \item \textbf{几何检查器}：在 isp\_tvalid 上数 tuser/tlast。
    tuser 到达即"结账"上一帧（tuser 恰 1、tlast 恰 720），
    结账不符打 \code{GEO\_ERR}；
  \item \textbf{检测检查器}：\code{t\_upd} 拍采样锁定结果，
    连续多帧坐标在框内则 \code{det\_ok} 计数；
  \item \textbf{UART 解码器}：一个 9 态状态机按 A5..5A 组帧，
    校验和 + valid + 坐标范围三项对齐才计 \code{u\_good}；
  \item \textbf{终判与超时}：帧数够、检测够、UART 够则 PASS；
    450ms 超时 FAIL 并打印各计数——\textbf{失败也要给出诊断数据}。
\end{itemize}

\begin{keybox}{为什么 TB 里要有 UART 解码器}
外设层（遥测）是唯一能"整链端到端"观测检测结果的通道——
屏幕上看到的是 OSD，日志里看到的是 UART。两个通道互为佐证，
也提前验证了第二册的遥测协议。判据 D 用的是同一个
\code{uart\_rx\_mon}，一个监视器两处复用。
\end{keybox}

% ============================================================
\section{迭代复盘（本册相关 5 轮）}

\begin{center}
\begin{tabular}{llp{5.0cm}l}
\toprule
\rowcolor{rowtint}\textbf{轮} & \textbf{现象} & \textbf{根因 / 处置} & \textbf{类型} \\
\midrule
1 & elaboration 58840 错误 & 门级网表 + TD 原语 + 缺拷贝；
  写两个行为替身 + 补 data\_96bit\_to\_128bit & 环境 \\
2 & I\_tready 接常量报错 & 端口是输出，命名反向 & 集成 \\
3 & 重复声明 & 先赋值后声明，隐式 wire 冲突 & RTL \\
7 & 整链 tuser 全灭 & 官方 demosaic 依赖真实 CSI 时序；
  列 TD 阶段项 + 检查点文档化 & 界碑 \\
8 & RAW10 生成器只出 1 帧 & heredoc 转义吞替换；
  改行级补丁修复（4 帧 921600 拍） & 脚本 \\
\bottomrule
\end{tabular}
\end{center}

\begin{warnbox}{轮 1 的启示：58840 个错误不是"改 58840 次"}
三个根因、两个替身、一次补拷贝——\textbf{错误数量的规模不代表工作量}
。遇到海量报错先聚类：58840 条里其实只有 3 类。这个习惯在
RTL 调试里到处适用。
\end{warnbox}

\begin{tipbox}{脚本类 bug 的通用教训（轮 8）}
生成器出错时不要用"更复杂的转义"去绕——直接\textbf{改行级补丁}
（读行、替换指定行、写回）。本仓库为此吃过三次亏
（Step2 的 fix 脚本改坏 tex、本轮 heredoc 吞替换），
结论写进方法论：\textbf{改文件用行级操作，别用嵌套转义。}
\end{tipbox}

% ============================================================
\section{上板检查点与已知差异（交接给上板阶段）}

\begin{center}
\begin{tabular}{ll}
\toprule
\rowcolor{rowtint}\textbf{现象（上板若遇到）} & \textbf{先查这里} \\
\midrule
有数据无帧同步 & demosaic tuser 链（本仿真不激活的同一问题） \\
时序不收敛 & divider\_gate 门级网表路径；例程本身未收敛为已知现象 \\
跟踪颜色漂移 & AWB 自适应（对策：锁曝光/切阈值组 8 组） \\
画面随舵机动作闪 & 舵机供电（必须独立 5V 共地，仿真无法暴露） \\
\bottomrule
\end{tabular}
\end{center}

\begin{okbox}{本册交付清单}
\begin{itemize}
  \item \code{system\_top.v}：集成顶层（elaboration 零错误）
  \item 两个行为替身 + filelist 的"仿真专用"清单
  \item 虚拟 RAW10 传感器生成器（固定种子可复现）
  \item \code{tb\_system\_integration}：四判据 + 超时诊断
  \item 分层证明材料与上板检查点清单
\end{itemize}
\end{okbox}
